% =====================================================================
%  preamble.tex —— 导言区（Lead 独占；译者只读，不得修改）
%  寺杣友秀《コンツェビッチによるウィッテン予想の解決》
%  （上智大学数学講究録 No.40, 1997）中文译注本 · XeLaTeX 编译
% =====================================================================
\documentclass[11pt,a4paper]{article}

\usepackage[UTF8]{ctex}
\usepackage{amsmath,amssymb,amsthm,mathtools}
\usepackage{mathrsfs}
\usepackage{bm}
\usepackage{graphicx}
\usepackage{tikz}
\usepackage{tikz-cd}
\usepackage[a4paper,left=2.5cm,right=2.5cm,top=2.7cm,bottom=2.8cm]{geometry}
\usepackage[dvipsnames]{xcolor}
\usepackage{enumitem}
\usepackage{booktabs}
\usepackage{array}
\usepackage{multicol}
\usepackage{fancyhdr}
\usepackage{caption}
\usepackage{titlesec}
\usepackage[hidelinks,unicode]{hyperref}
\hypersetup{pdftitle={Kontsevich 对 Witten 猜想的解决（中文译注本）},
  pdfauthor={寺杣友秀 著；中文译注：本工程}}

\captionsetup{font=small,labelsep=quad}
\setlength{\emergencystretch}{3em}
\allowdisplaybreaks[1]
\graphicspath{{figures/}}

% ---------- 节标题：原文「§1 …」→「§1 …」；小节「1.1 …」 ----------
\titleformat{\section}{\Large\bfseries}{\S\thesection}{0.8em}{}
\titlespacing*{\section}{0pt}{20pt}{12pt}
\titleformat{\subsection}{\large\bfseries}{\thesubsection}{0.8em}{}
\titlespacing*{\subsection}{0pt}{16pt}{8pt}
\titleformat{\subsubsection}{\normalsize\bfseries}{\thesubsubsection}{0.8em}{}

% ---------- 页眉页脚 ----------
\pagestyle{fancy}
\fancyhf{}
\fancyhead[L]{\small\itshape Witten's Intersection Theory of Moduli Spaces of Curves}
\fancyhead[R]{\small 寺杣友秀\quad\thepage}
\renewcommand{\headrulewidth}{0.4pt}

% ---------- 定理类环境（与原书编号规则一致，实测） ----------
% 原书编号实测：
%  * 定义 Definition 单独计数，按节编号（定义 1.1 …1.5；2.1 …；3.1,3.2）
%  * 定理 Theorem / 引理 Lemma / 命题 Proposition / 推论 Corollary 共用
%    一个按节编号的计数器（定理1.6 → 引理1.7 → 定理1.8；命题3.1 →
%    定理3.2 → 定理3.3 → 命题3.4 → 定理3.5 → 命题3.6 → 引理3.7 → 定理3.8）
%  * 例 Example 全文连续编号，不带节号（例 1、例 2）
%  * 注 Remark 无编号；问题 Problem 按节编号（原书未出现）
%  * 特例：原书第 16 页的「Proposition 1.8」与第 14 页的「Theorem 1.8」
%    重号（原书自身如此），译本用 \addtocounter{theorem}{-1} 在该处单点
%    复现，并在译注中说明，详见 chunks 中该处注释。
\theoremstyle{definition}
\newtheorem{definition}{定义}[section]
\newtheorem{theorem}{定理}[section]
\newtheorem{lemma}[theorem]{引理}
\newtheorem{proposition}[theorem]{命题}
\newtheorem{corollary}[theorem]{推论}
\newtheorem{example}{例}
\newtheorem*{remark}{注}
\newtheorem{problem}{问题}[section]
\renewcommand{\proofname}{证明}

% ---------- 常用宏 ----------
\newcommand{\R}{\mathbb{R}}
\newcommand{\C}{\mathbb{C}}
\newcommand{\Z}{\mathbb{Z}}
\newcommand{\N}{\mathbb{N}}
\newcommand{\Q}{\mathbb{Q}}
\newcommand{\PP}{\mathbb{P}}
\newcommand{\Spec}{\operatorname{Spec}}
\newcommand{\dd}{\mathrm{d}}
\newcommand{\id}{\mathrm{id}}
\newcommand{\rank}{\operatorname{rank}}
\newcommand{\Hom}{\operatorname{Hom}}
\newcommand{\Ext}{\operatorname{Ext}}
\newcommand{\Aut}{\operatorname{Aut}}
\newcommand{\Isom}{\operatorname{Isom}}
\newcommand{\Pic}{\operatorname{Pic}}
\newcommand{\Sch}{\operatorname{Sch}}
\newcommand{\Ob}{\operatorname{Ob}}
\newcommand{\End}{\operatorname{End}}
\newcommand{\Tr}{\operatorname{Tr}}
\newcommand{\tr}{\operatorname{tr}}
\newcommand{\Coeq}{\operatorname{Coeq}}
\newcommand{\coker}{\operatorname{coker}}
\newcommand{\vol}{\operatorname{vol}}
\newcommand{\et}{\text{\'et}}
% 模空间与模叠记号
\newcommand{\Mg}{M_{g}}
\newcommand{\Mbar}{\overline{\mathcal{M}}}
\newcommand{\MbarG}{\overline{\mathcal{M}}_{g}}
\newcommand{\MbarGN}{\overline{\mathcal{M}}_{g,n}}
\newcommand{\MbarN}[1]{\overline{\mathcal{M}}_{#1}}
\newcommand{\cM}{\mathcal{M}}
% 常用花体
\newcommand{\cA}{\mathcal{A}}\newcommand{\cB}{\mathcal{B}}\newcommand{\cC}{\mathcal{C}}
\newcommand{\cD}{\mathcal{D}}\newcommand{\cE}{\mathcal{E}}\newcommand{\cF}{\mathcal{F}}
\newcommand{\cG}{\mathcal{G}}\newcommand{\cH}{\mathcal{H}}\newcommand{\cI}{\mathcal{I}}
\newcommand{\cJ}{\mathcal{J}}\newcommand{\cK}{\mathcal{K}}\newcommand{\cL}{\mathcal{L}}
\newcommand{\cN}{\mathcal{N}}\newcommand{\cO}{\mathcal{O}}\newcommand{\cP}{\mathcal{P}}
\newcommand{\cQ}{\mathcal{Q}}\newcommand{\cR}{\mathcal{R}}\newcommand{\cS}{\mathcal{S}}
\newcommand{\cT}{\mathcal{T}}\newcommand{\cU}{\mathcal{U}}\newcommand{\cV}{\mathcal{V}}
\newcommand{\cW}{\mathcal{W}}\newcommand{\cX}{\mathcal{X}}\newcommand{\cY}{\mathcal{Y}}
\newcommand{\cZ}{\mathcal{Z}}

% 译者注（就地批注）：命令形式，不要写成环境
% 用法：\transnote{注文}
\newcommand{\transnote}[1]{\par\medskip\noindent{\small\color{Sepia}\textbf{【译者注】}#1\par}\medskip}

\setlength{\parindent}{2em}
\setcounter{tocdepth}{2}
